\documentclass[11pt]{article}
\usepackage[T1]{fontenc}
\usepackage{lmodern}
\usepackage[margin=1.05in]{geometry}
\usepackage{amsmath,amssymb,amsthm}
\usepackage{booktabs}
\usepackage{microtype}
\microtypesetup{expansion=false}
\usepackage{xcolor}
\definecolor{linknavy}{RGB}{25,55,125}
\usepackage[colorlinks=true,linkcolor=linknavy,citecolor=linknavy,urlcolor=linknavy]{hyperref}
\hypersetup{
  pdfauthor={Felipe Santibáñez-Leal},
  pdftitle={Tensor Factorization of Extension Ideals under Gluing and Huneke--Wiegand Counterexamples with Many Generators},
  pdfsubject={Gluings of numerical semigroups, torsion in tensor products, the Huneke-Wiegand conjecture},
  pdfkeywords={Huneke-Wiegand conjecture, numerical semigroups, gluing, semigroup rings, torsion, tensor products, relative ideals}
}
\usepackage{orcidlink}

\newtheorem{theorem}{Theorem}[section]
\newtheorem{proposition}[theorem]{Proposition}
\newtheorem{corollary}[theorem]{Corollary}
\newtheorem{lemma}[theorem]{Lemma}
\theoremstyle{definition}
\newtheorem{definition}[theorem]{Definition}
\newtheorem{example}[theorem]{Example}
\theoremstyle{remark}
\newtheorem{remark}[theorem]{Remark}

\newcommand{\N}{\mathbb{N}}
\newcommand{\Z}{\mathbb{Z}}
\newcommand{\kk}{\Bbbk}
\newcommand{\Ap}{\operatorname{Ap}}
\newcommand{\tor}{\tau}
\newcommand{\len}{\ell}

% Page-1 header block (conventions/manuscript-header-standard.md). The four document
% macros are set by the publishing tool when the version DOI is reserved.
\newcommand{\doctype}{Preprint}
\newcommand{\docversion}{v0.01}
\newcommand{\docdate}{30 October 2026}
\newcommand{\versiondoi}{10.5281/zenodo.23821498}
\newcommand{\conceptdoi}{10.5281/zenodo.23821497}
\newcommand{\docrepo}{github.com/fsantibanezleal/CAOS\_RESEARCH}

\newcommand{\headerblock}{%
\begin{center}
\fbox{\parbox{0.94\textwidth}{\small\raggedright
\textbf{\doctype} \quad$\cdot$\quad \docversion, \docdate \quad$\cdot$\quad Licence CC BY 4.0\\[2pt]
CAOS Research open-problems programme \quad$\cdot$\quad problem \texttt{huneke-wiegand}\\[2pt]
Version DOI \href{https://doi.org/\versiondoi}{\texttt{\versiondoi}}\\[2pt]
Concept DOI (always latest) \href{https://doi.org/\conceptdoi}{\texttt{\conceptdoi}}\\[2pt]
Not peer reviewed. Source code, data and computational records: \href{https://\docrepo}{\texttt{\docrepo}}
}}
\end{center}
\vspace{4pt}}

\title{Tensor Factorization of Extension Ideals under Gluing\\
and Huneke--Wiegand Counterexamples with Many Generators}
\author{Felipe Santib\'a\~nez-Leal\,\orcidlink{0000-0002-0150-3246}\\
\small CAOS Research programme, \texttt{github.com/fsantibanezleal/CAOS\_RESEARCH}}
\date{\docdate}

\begin{document}
\maketitle
\headerblock

\begin{abstract}
Let $\Gamma=a_1\Gamma_1+a_2\Gamma_2$ be a gluing of numerical semigroups, $\kk$ a field, and
$A=a_1A_1+a_2A_2$ the extension of relative ideals $A_i$ of $\Gamma_i$. With $y$ acting on
$\kk[\Gamma_1]$ as $t^{a_2}$ and on $\kk[\Gamma_2]$ as $t^{a_1}$, we prove that
$I_A\otimes_{\kk[\Gamma]} I_{A^*}\cong T_1\otimes_{\kk[y]}T_2$, where
$T_i=I_{A_i}\otimes I_{A_i^*}$. Consequently the torsion lengths satisfy
\[
 \len(\tor(I_A\otimes I_{A^*}))=a_1\len(\tor(T_1))+a_2\len(\tor(T_2))+\sum_{j\ge1}b_1(j)\,b_2(j),
\]
where $b_i(j)$ counts the Jordan blocks of size at least $j$ of $y$ on $\tor(T_i)$, and the numbers of
minimal generators multiply. Hence $A$ has the Huneke--Wiegand property if and only if $A_1$ or $A_2$
is a non-principal ideal with that property. This re-proves the two gluing propositions of
Garc\'{\i}a-S\'anchez and Leamer and settles the case they leave open: when both factors are
non-principal and the tensor product of each with its dual is torsion-free, so is the tensor product for
the extension. Since the Huneke--Wiegand conjecture is now known to be false, this produces
counterexample ideals with $2^r$ minimal generators over symmetric numerical semigroup rings. We give
two explicit four-generated examples, over symmetric semigroups with Frobenius numbers $23645$ and
$53599$, built from Pham's counterexample and from an infinite family of counterexamples. Exact
computations by two independent routes, the definition of the tensor product and Leamer's partition
criterion, confirm the length law on 200 seeded gluings.
\end{abstract}

\section{Introduction}

Let $R$ be a one-dimensional Gorenstein local domain. Huneke and Wiegand asked whether a finitely
generated torsion-free $R$-module $M$ with $M\otimes_R\operatorname{Hom}_R(M,R)$ torsion-free must be
free \cite{HunekeWiegand1994}. For numerical semigroup rings and monomial ideals, Leamer reduced the
question to arithmetic of relative ideals \cite{Leamer2012}, and Garc\'{\i}a-S\'anchez and Leamer
proved it for monomial ideals of free numerical semigroup rings and for two-generated monomial ideals
over complete intersections \cite{GarciaSanchezLeamer2013}. Their method rests on gluings
$\Gamma=a_1\Gamma_1+a_2\Gamma_2$: for an extension $a_1A_1+a_2A_2$ they show that it inherits the
Huneke--Wiegand property from a non-principal factor (their Proposition~9), and that
$a_1A_1+a_2\Gamma_2$ has the property exactly when $A_1$ does (their Proposition~10). In their final
section they ask whether the property for ideals with more generators is invariant under gluing, and
whether their Proposition~9 extends to rings of the form
$(R_1\otimes_\kk R_2)/(f\otimes1-1\otimes g)$.

The conjecture is now known to be false in dimension one. Christensen, Gerko and Iyengar gave a
non-principal rigid ideal in a one-dimensional Gorenstein domain \cite{ChristensenGerkoIyengar2026},
and Pham gave a two-generated monomial counterexample over a symmetric numerical semigroup of
multiplicity $56$ and Frobenius number $181$ \cite{Pham2026}. An infinite family of two-generated
counterexamples of multiplicity $24p$, $p\ge4$, is constructed in \cite{SantibanezLeal2026}.

In the presence of counterexamples, the case of an extension whose two factors are both non-principal
and both fail the Huneke--Wiegand property is no longer vacuous, and it is not covered by the
propositions above. We prove a tensor factorization that decides every extension.

\begin{theorem}\label{thm:main}
Let $\Gamma=a_1\Gamma_1+a_2\Gamma_2$ be a gluing, $A_i$ relative ideals of $\Gamma_i$, and
$A=a_1A_1+a_2A_2$. Write $T=I_A\otimes I_{A^*}$ and $T_i=I_{A_i}\otimes I_{A_i^*}$, and let $y$ act on
$\kk[\Gamma_1]$ as $t^{a_2}$ and on $\kk[\Gamma_2]$ as $t^{a_1}$. Then
\begin{enumerate}
\item $T\cong T_1\otimes_{\kk[y]}T_2$;
\item $\len(\tor(T))=a_1\len(\tor(T_1))+a_2\len(\tor(T_2))+\sum_{j\ge1}b_1(j)b_2(j)$, where $b_i(j)$ is
the number of Jordan blocks of size at least $j$ of $y$ acting on $\tor(T_i)$;
\item $\mu(A)=\mu(A_1)\mu(A_2)$;
\item $A$ has the Huneke--Wiegand property if and only if $A_1$ or $A_2$ is a non-principal ideal with
that property.
\end{enumerate}
\end{theorem}

\begin{corollary}\label{cor:many}
If $A_1$ and $A_2$ are two-generated counterexample ideals over symmetric numerical semigroups
$\Gamma_1$ and $\Gamma_2$, and $\Gamma=a_1\Gamma_1+a_2\Gamma_2$ is a gluing, then $a_1A_1+a_2A_2$ is a
four-generated counterexample ideal over the symmetric semigroup $\Gamma$. Iterating, for every $r\ge1$
there are counterexample ideals with $2^r$ minimal generators over symmetric numerical semigroup
rings.
\end{corollary}

Section~\ref{sec:prelim} fixes notation. Section~\ref{sec:factor} proves part~(1), and
Section~\ref{sec:length} parts (2) to (4). Section~\ref{sec:consequences} derives the propositions of
\cite{GarciaSanchezLeamer2013}, Corollary~\ref{cor:many}, two explicit examples and a partial answer to
their invariance question. Section~\ref{sec:computation} reports exact computations, and
Section~\ref{sec:scope} states the limits of the results.

\section{Preliminaries}\label{sec:prelim}

A numerical semigroup $\Lambda\subseteq\N$ is a submonoid with finite complement; $F(\Lambda)$ is its
Frobenius number. A relative ideal is a non-empty set $B\subseteq\Z$, bounded below, with
$B+\Lambda\subseteq B$. Its dual is $B^*=\Lambda-B=\{z:z+B\subseteq\Lambda\}$, and $B$ is principal if
$B=x+\Lambda$. The fractional ideal $I_B$ of $\kk[\Lambda]$ is spanned by the monomials $t^b$, $b\in B$,
and $I_{B^*}\cong\operatorname{Hom}(I_B,\kk[\Lambda])$ \cite[Remark~1.7]{Leamer2012}. Following
\cite{GarciaSanchezLeamer2013}, a non-principal $B$ has the \emph{Huneke--Wiegand property} if
$I_B\otimes I_{B^*}$ has non-zero torsion; by \cite[Theorem~1.4]{Leamer2012} this holds iff the
minimal generators of $B$ split into two non-empty sets generating $P$ and $Q$ with
$(P+B^*)\cap(Q+B^*)\ne(P\cap Q)+B^*$. A \emph{counterexample ideal} is a non-principal $B$ without the
property. We write $\mu(B)$ for the number of minimal generators and
$\Ap(B,c)=\{b\in B:b-c\notin B\}$ for $c\in\Lambda$, $c>0$.

$\Gamma$ is the \emph{gluing} of $\Gamma_1$ and $\Gamma_2$ if $\Gamma=a_1\Gamma_1+a_2\Gamma_2$ with
$a_1\in\Gamma_2$, $a_2\in\Gamma_1$ and $\gcd(a_1,a_2)=1$; then $F(\Gamma)=a_1F(\Gamma_1)+a_2F(\Gamma_2)+a_1a_2$,
and $\Gamma$ is symmetric when $\Gamma_1$ and $\Gamma_2$ are \cite{Delorme1976,RosalesGarciaSanchez2009}.
The \emph{extension} of $A_i\subseteq\Gamma_i$ is $A=a_1A_1+a_2A_2$; it is a relative ideal of $\Gamma$,
generated by $a_1x+a_2z$ over generators $x$ of $A_1$ and $z$ of $A_2$, and
$A^*=a_1A_1^*+a_2A_2^*$ \cite[Lemma~8]{GarciaSanchezLeamer2013}. Throughout, $R_i=\kk[\Gamma_i]$,
$R=\kk[\Gamma]$, $S=\kk[y]$ with $y\mapsto t^{a_2}$ in $R_1$ and $y\mapsto t^{a_1}$ in $R_2$. All rings
and modules are graded by the exponent; for graded modules of finite length, $\len=\dim_\kk$.

\section{Extensions are tensor products}\label{sec:factor}

\begin{lemma}\label{lem:apery}
Let $B$ be a relative ideal of $\Lambda$ and $c\in\Lambda$, $c>0$. Then $I_B$ is a free
$\kk[t^c]$-module with basis $\{t^b:b\in\Ap(B,c)\}$, and $\Ap(B,c)$ has exactly one element in each
residue class modulo $c$.
\end{lemma}

\begin{proof}
Each residue class modulo $c$ meets $B$ in a non-empty set that is bounded below and closed under
adding $c$. With $b_r$ its least element, the class is $b_r+c\N$, so $I_B=\bigoplus_r t^{b_r}\kk[t^c]$.
\end{proof}

\begin{proposition}\label{prop:iso}
The assignment $t^x\otimes t^z\mapsto t^{a_1x+a_2z}$ defines an isomorphism of
$R_1\otimes_SR_2$-modules $I_{A_1}\otimes_SI_{A_2}\to I_A$. For $A_i=\Gamma_i$ it is a ring isomorphism
$R_1\otimes_SR_2\cong R$, and $I_{A^*}\cong I_{A_1^*}\otimes_SI_{A_2^*}$.
\end{proposition}

\begin{proof}
The map is $S$-balanced, since $(t^{a_2}t^x)\otimes t^z$ and $t^x\otimes(t^{a_1}t^z)$ both go to
$t^{a_1x+a_1a_2+a_2z}$, and it is compatible with both actions. By Lemma~\ref{lem:apery} the source is
$S$-free on $t^x\otimes t^z$ with $x\in\Ap(A_1,a_2)$ and $z\in\Ap(A_2,a_1)$, so it has the $\kk$-basis
$y^n(t^x\otimes t^z)$, sent to $t^{a_1x+a_2z+na_1a_2}$. Every element of $A$ is $a_1x'+a_2z'$ with
$x'=x+ia_2$ and $z'=z+ja_1$ for Ap\'ery elements $x,z$ and $i,j\ge0$, hence equals
$a_1x+a_2z+(i+j)a_1a_2$; so the map is onto. If $a_1x+a_2z+na_1a_2=a_1x'+a_2z'+n'a_1a_2$, reduction
modulo $a_1$ gives $z\equiv z'$, hence $z=z'$ by Lemma~\ref{lem:apery}; likewise $x=x'$, and then
$n=n'$. Distinct basis vectors therefore go to distinct monomials. The last statement follows from
$A^*=a_1A_1^*+a_2A_2^*$.
\end{proof}

\begin{proposition}\label{prop:tensor}
$T\cong T_1\otimes_ST_2$ as modules over $R\cong R_1\otimes_SR_2$.
\end{proposition}

\begin{proof}
For commutative $S$-algebras $R_i$ and $R_i$-modules $M_i,N_i$ the assignment
$(m_1\otimes m_2)\otimes(n_1\otimes n_2)\mapsto(m_1\otimes n_1)\otimes(m_2\otimes n_2)$ gives an
isomorphism
$(M_1\otimes_SM_2)\otimes_{R_1\otimes_SR_2}(N_1\otimes_SN_2)\cong(M_1\otimes_{R_1}N_1)\otimes_S(M_2\otimes_{R_2}N_2)$:
both sides are the quotient of $M_1\otimes_SM_2\otimes_SN_1\otimes_SN_2$ by the relations
$r_im_i\otimes n_i=m_i\otimes r_in_i$, and the inverse assignment is well defined on them. Apply it with
$M_i=I_{A_i}$, $N_i=I_{A_i^*}$ and Proposition~\ref{prop:iso}.
\end{proof}

\section{The torsion-length law}\label{sec:length}

\begin{lemma}\label{lem:torsion}
$S$ embeds in $R$ by $y\mapsto t^{a_1a_2}$, $R$ is a finite $S$-module, and $R_1$ and $R_2$ are free over
$S$ of ranks $a_2$ and $a_1$. For a finitely generated module over $R$, $R_1$ or $R_2$, the torsion
submodule equals the $S$-torsion submodule.
\end{lemma}

\begin{proof}
The ranks come from Lemma~\ref{lem:apery}. A non-zero element of $S$ stays non-zero in each domain, so
$S$-torsion is torsion. Conversely, if $rm=0$ with $r\ne0$, take a monic relation
$r^d+s_{d-1}r^{d-1}+\cdots+s_0=0$ over $S$ of least degree; then $s_0\ne0$ and
$s_0m=-(r^{d-1}+\cdots+s_1)rm=0$.
\end{proof}

\begin{proof}[Proof of Theorem~\ref{thm:main}]
Part~(1) is Proposition~\ref{prop:tensor}. For (2), each $T_i$ is finitely generated over $S$ and has
rank one over $R_i$, so over the graded principal ideal domain $S$ we have
$T_1\cong S^{a_2}\oplus\tor_1$ and $T_2\cong S^{a_1}\oplus\tor_2$ with $\tor_i=\tor(T_i)$ by
Lemma~\ref{lem:torsion}. Then
\[
T\cong S^{a_1a_2}\oplus\tor_1^{\,a_1}\oplus\tor_2^{\,a_2}\oplus(\tor_1\otimes_S\tor_2).
\]
The last three summands are $S$-torsion and the first is free, so their sum is $\tor(T)$. Each $\tor_i$
has finite length and $y$ raises degree, so $y$ is nilpotent on it:
$\tor_1\cong\bigoplus S/(y^{e_i})$ and $\tor_2\cong\bigoplus S/(y^{f_j})$ up to shifts. Since
$S/(y^e)\otimes_SS/(y^f)\cong S/(y^{\min(e,f)})$, the length of $\tor_1\otimes_S\tor_2$ is
$\sum_{i,j}\min(e_i,f_j)=\sum_{j\ge1}b_1(j)b_2(j)$.

For (3), let $\mathfrak m$ be the graded maximal ideal. Then $R/\mathfrak m=\kk=(R_1/\mathfrak m_1)\otimes_S(R_2/\mathfrak m_2)$
with $y$ acting by zero, and the isomorphism of Proposition~\ref{prop:tensor} applied to $N_i=\kk$ gives
$I_A\otimes_R\kk\cong(I_{A_1}\otimes_{R_1}\kk)\otimes_\kk(I_{A_2}\otimes_{R_2}\kk)$, of dimension
$\mu(A_1)\mu(A_2)$.

For (4), a principal factor has $T_i\cong R_i$, which is torsion-free, and a non-principal factor has
$\tor(T_i)\ne0$ exactly when it has the Huneke--Wiegand property. Since $a_1,a_2\ge1$, part~(2) gives
$\tor(T)\ne0$ iff $\tor(T_1)\ne0$ or $\tor(T_2)\ne0$, and part~(3) gives that $A$ is non-principal iff a
factor is.
\end{proof}

\section{Consequences}\label{sec:consequences}

\begin{corollary}[{\cite[Propositions 9 and 10]{GarciaSanchezLeamer2013}}]
If $A_1$ is a non-principal ideal with the Huneke--Wiegand property, every extension $a_1A_1+a_2B$ has
it. The extension $a_1A_1+a_2\Gamma_2$ has it if and only if $A_1$ has.
\end{corollary}

\begin{corollary}\label{cor:open}
If $A_1$ and $A_2$ are both counterexample ideals, then $a_1A_1+a_2A_2$ is a counterexample ideal with
$\mu(A_1)\mu(A_2)$ minimal generators.
\end{corollary}

Corollary~\ref{cor:many} follows from Corollary~\ref{cor:open} and the symmetry of gluings of symmetric
semigroups. For the simple gluings $m\N+q\Gamma$ with $\gcd(q,m)=1$ and $m$ the multiplicity of
$\Gamma$, the case $\Gamma_2=\N$, $A_2=\N$ of Theorem~\ref{thm:main}, already given by
\cite[Proposition~10]{GarciaSanchezLeamer2013}, turns the family of \cite{SantibanezLeal2026} into a
two-parameter family of symmetric counterexamples of multiplicity $24p$ and embedding dimension $11p$,
for every $p\ge4$ and every $q\ge1$ coprime to $24p$.

\begin{example}\label{ex:pham}
Let $\Gamma_P$ be Pham's semigroup, minimally generated by $56$, $57$, $58$, $63$, $64$, the integers from $70$ to $83$, and $87$, $89$,
$90$, $93$, $95$, $96$, $97$ (multiplicity $56$, Frobenius number $181$), with the counterexample ideal $(0,14)$ \cite{Pham2026}. For
the gluing $\Gamma=56\Gamma_P+57\Gamma_P$, the extension $56(0,14)+57(0,14)$ is the relative ideal with
minimal generators $0,784,798,1582$. $\Gamma$ is symmetric with Frobenius number
$56\cdot181+57\cdot181+56\cdot57=23645$, and the extension is a four-generated counterexample.
\end{example}

\begin{example}\label{ex:mixed}
Let $C_4$ be the member $p=4$ of the family of \cite{SantibanezLeal2026}, of multiplicity $96$ and
Frobenius number $311$, with counterexample ideal $(0,24)$. For $\Gamma=96\Gamma_P+89C_4$, the extension
$96(0,14)+89(0,24)$ has minimal generators $0,1344,2136,3480$, and $\Gamma$ is symmetric with Frobenius
number $96\cdot181+89\cdot311+96\cdot89=53599$.
\end{example}

\begin{corollary}
If every non-principal relative ideal of $\Gamma_1$ and of $\Gamma_2$ has the Huneke--Wiegand property,
then so does every non-principal extension ideal of $a_1\Gamma_1+a_2\Gamma_2$.
\end{corollary}

This answers the invariance question of \cite[Section~4]{GarciaSanchezLeamer2013} for extension ideals.
Relative ideals of $\Gamma$ that are not extensions are not covered.

\begin{remark}
For monomial ideals, the module $(I_{A_1}\otimes_\kk I_{A_2})/(t^{a_2}\otimes1-1\otimes t^{a_1})$ of the
ring-level question in \cite[Section~4]{GarciaSanchezLeamer2013} is $I_{A_1}\otimes_SI_{A_2}=I_A$, so
Theorem~\ref{thm:main} answers that question for monomial ideals. A statement for arbitrary modules would
need the dual of $M_1\otimes_SM_2$ to factor over $S$, which is not established here.
\end{remark}

\section{Computational checks}\label{sec:computation}

The statements above are proved; the computations below test them and their transcription. All
arithmetic is exact.

\emph{Two routes.} The torsion of $I_B\otimes I_{B^*}$ is computed from the definition: in degree $d$
the module is spanned by $e_x\otimes e_z$ with $x\in B$, $z\in B^*$, $x+z=d$, modulo
$e_{x+g}\otimes e_z=e_x\otimes e_{z+g}$ for generators $g$, so its dimension is the number of connected
components of a finite graph and the torsion in degree $d$ has one dimension fewer. Every pair joins a
pair $(x_i,d-x_i)$ with $x_i$ a generator of $B$, and two such pairs join through $w=\max_ix_i+F(\Lambda)+1$
once $d-w$ exceeds the conductor of $B^*$, which bounds the degrees to scan. Independently, the
Huneke--Wiegand property is decided by the partition criterion of \cite{Leamer2012}.

\emph{Seeded family.} For 200 gluings with random factors of Frobenius number at most $10$, gluing
parameters at most $7$ and ideals with at most three generators, the directly computed torsion length of
the extension equals part~(2) of Theorem~\ref{thm:main} computed from the factors alone, including the
Jordan types of $y$; in 33 of the cases both factors have torsion, so the cross term is non-zero. The
partition criterion agrees with part~(4), and part~(3) and the Hilbert series identity
$H_{I_A}(z)=H_{I_{A_1}}(z^{a_1})H_{I_{A_2}}(z^{a_2})(1-z^{a_1a_2})$ hold in every case.

\emph{Examples and controls.} The partition criterion confirms Examples~\ref{ex:pham} and~\ref{ex:mixed},
and that $(0,14)$ and $(0,24)$ are counterexample ideals of $\Gamma_P$ and $C_4$. Replacing the second
factor of Example~\ref{ex:pham} by the ideal $(0,1)$ of $\Gamma_P$, which has the property, yields an
extension with the property. For 30 random pairs with $a_1\notin\Gamma_2$ the Hilbert series identity
fails, as it must without the gluing hypothesis. A second implementation sharing no code with the first
recomputes Example~\ref{ex:pham} and the torsion lengths of 25 of the 200 cases, with full agreement. The
code and data are in the repository named on the first page.

\section{Scope}\label{sec:scope}

The results concern monomial (relative) ideals of numerical semigroup rings over any field, graded or
completed. Theorem~\ref{thm:main} decides extension ideals of a gluing only; whether the gluing of two
semigroups whose non-principal ideals all have the Huneke--Wiegand property can carry a counterexample
ideal that is not an extension remains open. The sources consulted for the factorization and the length
law were \cite{GarciaSanchezLeamer2013,GotoTakahashiTaniguchiTruong2013,Leamer2012,Roy2018}; none of them
states these results, which is not a proof that they are new.

\section*{Acknowledgments}

The author thanks Son Pham for making the first numerical-semigroup counterexample and its verification
materials public. Neither he nor the other authors cited is an author of this paper, and this
acknowledgment does not imply endorsement of the present results.

\medskip
\noindent\textbf{Declaration of generative AI use.} Generative AI tools were used in the
preparation of this paper. The author reviewed and verified all content and takes full
responsibility for the work.

{\footnotesize
\begin{thebibliography}{99}
\setlength{\itemsep}{0.25em}
\setlength{\parsep}{0pt}

\bibitem{ChristensenGerkoIyengar2026}
L. W. Christensen, A. Gerko and S. B. Iyengar,
\emph{Some non-principal rigid ideals in Gorenstein domains of dimension one},
preprint (2026), \href{https://arxiv.org/abs/2608.21666}{arXiv:2608.21666}.

\bibitem{Delorme1976}
C. Delorme,
\emph{Sous-mono\"ides d'intersection compl\`ete de $\N$},
Ann. Sci. \'Ec. Norm. Sup\'er. (4) 9 (1976), 145--154.

\bibitem{GarciaSanchezLeamer2013}
P. A. Garc\'{\i}a-S\'anchez and M. J. Leamer,
\emph{Huneke--Wiegand conjecture for complete intersection numerical semigroup rings},
J. Algebra 391 (2013), 114--124; \href{https://arxiv.org/abs/1211.4554}{arXiv:1211.4554}.

\bibitem{GotoTakahashiTaniguchiTruong2013}
S. Goto, R. Takahashi, N. Taniguchi and H. Le Truong,
\emph{Huneke--Wiegand conjecture and change of rings},
preprint (2013), \href{https://arxiv.org/abs/1305.4238}{arXiv:1305.4238}.

\bibitem{HunekeWiegand1994}
C. Huneke and R. Wiegand,
\emph{Tensor products of modules and the rigidity of Tor},
Math. Ann. 299 (1994), 449--476.
\href{https://doi.org/10.1007/BF01459793}{doi:10.1007/BF01459793}.

\bibitem{Leamer2012}
M. J. Leamer,
\emph{Torsion and tensor products over domains and specializations to semigroup rings},
preprint (2012), \href{https://arxiv.org/abs/1211.2896}{arXiv:1211.2896}.

\bibitem{Pham2026}
S. Pham,
\emph{A counterexample to the Huneke--Wiegand conjecture},
preprint (2026), \href{https://arxiv.org/abs/2609.07615}{arXiv:2609.07615}.

\bibitem{RosalesGarciaSanchez2009}
J. C. Rosales and P. A. Garc\'{\i}a-S\'anchez,
\emph{Numerical Semigroups}, Developments in Mathematics 20, Springer, New York, 2009.

\bibitem{Roy2018}
R. Roy,
\emph{Graded AR sequences and the Huneke--Wiegand conjecture},
preprint (2018), \href{https://arxiv.org/abs/1808.06600}{arXiv:1808.06600}.

\bibitem{SantibanezLeal2026}
F. Santib\'a\~nez-Leal,
\emph{Frobenius minimality, minimum-layer uniqueness, and an infinite family of
Huneke--Wiegand counterexamples}, CAOS Research preprint (2026).
\href{https://doi.org/10.5281/zenodo.21763582}{doi:10.5281/zenodo.21763582}.

\end{thebibliography}
}

\end{document}
